\documentclass[12pt,a4paper]{article}
\usepackage[utf8]{inputenc}
\usepackage[T1]{fontenc}
\usepackage[polish]{babel}
\usepackage{amsmath}
\usepackage{amssymb}
\usepackage{geometry}
\usepackage{booktabs}
\usepackage{array}
\usepackage{xcolor}
\usepackage{hyperref}
\usepackage{tikz}
\usepackage{pgfplots}
\pgfplotsset{compat=1.18}
\usepackage{siunitx}
\sisetup{locale=DE}
\usepackage{mdframed}
\usepackage{enumitem}
\usepackage{float}

\geometry{margin=2.5cm}

\hypersetup{
    colorlinks=true,
    linkcolor=blue!70!black,
    urlcolor=blue!70!black,
}

% ---- custom environments ----
\newmdenv[
  backgroundcolor=blue!6,
  linecolor=blue!50,
  linewidth=1.2pt,
  roundcorner=5pt,
  skipabove=6pt,
  skipbelow=6pt,
]{wzorbox}

\newmdenv[
  backgroundcolor=green!6,
  linecolor=green!50,
  linewidth=1.2pt,
  roundcorner=5pt,
  skipabove=6pt,
  skipbelow=6pt,
]{objasnbox}

% ---- helpers ----
\newcommand{\jednostka}[1]{\textbf{#1}}
\newcommand{\oznaczenie}[2]{%
  \item $#1$ -- #2%
}

% ====================================================================
\title{\textbf{Fotonika -- opracowanie zagadnień kolokwialnych}\\[4pt]
       \large Na podstawie pliku \texttt{wymagania.md}}
\author{}
\date{\today}
% ====================================================================

\begin{document}
\maketitle
\tableofcontents
\newpage

% ====================================================================
\section{Fotometria}
% ====================================================================

\subsection{Podstawowe pojęcia fotometryczne}

Fotometria zajmuje się pomiarami promieniowania elektromagnetycznego z uwzględnieniem
wrażliwości ludzkiego oka.  Trzy kluczowe wielkości to:
\textbf{światłość}, \textbf{strumień świetlny} i \textbf{natężenie oświetlenia}.

% ------- Światłość --------------------------------------------------
\subsubsection{Światłość $I$}

Światłość opisuje, ile strumienia świetlnego źródło emituje w danym kierunku
na jednostkę kąta bryłowego.

\begin{wzorbox}
\begin{equation}
  I = \frac{d\Phi}{d\Omega}
  \label{eq:swiatlosc}
\end{equation}
\end{wzorbox}

\begin{objasnbox}
\textbf{Objaśnienie symboli:}
\begin{itemize}[leftmargin=2em]
  \oznaczenie{I}{Światłość \quad [\textbf{cd}] -- kandela (jednostka podstawowa SI)}
  \oznaczenie{d\Phi}{nieskończenie mały strumień świetlny emitowany w stożek
                    \quad [\textbf{lm}] -- lumen}
  \oznaczenie{d\Omega}{nieskończenie mały kąt bryłowy,
                    $\Omega = A/r^2$ dla powierzchni $A$ w odległości $r$
                    \quad [\textbf{sr}] -- steradian}
\end{itemize}
\end{objasnbox}

Dla źródła \emph{izotropowego} (emitującego równomiernie we wszystkich kierunkach)
całkowity kąt bryłowy to $\Omega = 4\pi\,\si{sr}$, więc:
\begin{equation}
  I = \frac{\Phi}{4\pi}
  \label{eq:swiatlosc_izotropowe}
\end{equation}

% ------- Strumień świetlny ------------------------------------------
\subsubsection{Strumień świetlny $\Phi$}

Strumień świetlny jest miarą całkowitej mocy promieniowania odczuwanej
przez ludzkie oko -- uwzględnia krzywą czułości widmowej oka $V(\lambda)$.

\begin{wzorbox}
\begin{equation}
  \Phi = K_m \int_{0}^{\infty} \Phi_e(\lambda)\, V(\lambda)\, d\lambda
  \label{eq:strumien}
\end{equation}
\end{wzorbox}

\begin{objasnbox}
\textbf{Objaśnienie symboli:}
\begin{itemize}[leftmargin=2em]
  \oznaczenie{\Phi}{Strumień świetlny \quad [\textbf{lm}] -- lumen}
  \oznaczenie{K_m}{Maksymalna skuteczność świetlna $\approx 683\,\si{lm/W}$ przy
                   $\lambda = 555\,\si{nm}$}
  \oznaczenie{\Phi_e(\lambda)}{Spektralna gęstość mocy (strumień energetyczny)
                   \quad [\textbf{W/nm}]}
  \oznaczenie{V(\lambda)}{Krzywa widzialności oka (bezwymiarowa, $0\leq V \leq 1$)}
  \oznaczenie{\lambda}{Długość fali \quad [\textbf{m}] lub praktycznie [\textbf{nm}]}
\end{itemize}
\end{objasnbox}

Upraszczając: dla źródła o mocy promieniowania $P$ i skuteczności świetlnej $\eta$:
\begin{equation}
  \Phi = \eta \cdot P
\end{equation}

% ------- Natężenie oświetlenia --------------------------------------
\subsubsection{Natężenie oświetlenia $E$}

Natężenie oświetlenia (iluminancja) określa, ile strumienia świetlnego
pada na jednostkę powierzchni oświetlanego obiektu.

\begin{wzorbox}
\begin{equation}
  E = \frac{d\Phi}{dA}
  \label{eq:natezenie}
\end{equation}
\end{wzorbox}

\begin{objasnbox}
\textbf{Objaśnienie symboli:}
\begin{itemize}[leftmargin=2em]
  \oznaczenie{E}{Natężenie oświetlenia \quad [\textbf{lx}] -- luks $= \si{lm/m^2}$}
  \oznaczenie{d\Phi}{Strumień świetlny padający na element powierzchni \quad [\textbf{lm}]}
  \oznaczenie{dA}{Element powierzchni \quad [\textbf{m}$^2$]}
\end{itemize}
\end{objasnbox}

% ------- Relacje ----------------------------------------------------
\subsubsection{Relacje między wielkościami fotometrycznymi}

Dla \textbf{punktowego źródła światła} o światłości $I$ natężenie oświetlenia
powierzchni prostopadłej do promienia, w odległości $r$, wynosi:

\begin{wzorbox}
\begin{equation}
  E = \frac{I}{r^2}
  \label{eq:prawo_odwrotnych_kwadratow}
\end{equation}
\end{wzorbox}

Jest to \textbf{prawo odwrotnych kwadratów} -- oświetlenie maleje z kwadratem
odległości od źródła.

Gdy oświetlana powierzchnia jest nachylona pod kątem $\alpha$ do promienia
(kąt padania $\theta$ do normalnej powierzchni $\theta = 90^\circ - \alpha$):

\begin{wzorbox}
\begin{equation}
  E = \frac{I \cos\theta}{r^2}
  \label{eq:lambert}
\end{equation}
\end{wzorbox}

\begin{objasnbox}
\textbf{Objaśnienie symboli:}
\begin{itemize}[leftmargin=2em]
  \oznaczenie{E}{Natężenie oświetlenia \quad [\textbf{lx}]}
  \oznaczenie{I}{Światłość źródła \quad [\textbf{cd}]}
  \oznaczenie{r}{Odległość od źródła do powierzchni \quad [\textbf{m}]}
  \oznaczenie{\theta}{Kąt między promieniem a normalną do powierzchni
                      \quad [\textbf{rad}] lub [\textbf{°}]}
\end{itemize}
\end{objasnbox}

\paragraph{Zestawienie jednostek fotometrycznych:}

\begin{table}[H]
\centering
\renewcommand{\arraystretch}{1.4}
\begin{tabular}{lll}
\toprule
\textbf{Wielkość} & \textbf{Symbol} & \textbf{Jednostka SI} \\
\midrule
Światłość        & $I$      & kandela [\textbf{cd}]   \\
Strumień świetlny & $\Phi$  & lumen [\textbf{lm}] $=$ cd$\cdot$sr \\
Natężenie oświetlenia & $E$ & luks [\textbf{lx}] $=$ lm/m$^2$ \\
Luminancja       & $L$      & [\textbf{cd/m}$^2$]     \\
\bottomrule
\end{tabular}
\caption{Podstawowe wielkości fotometryczne}
\end{table}

\newpage
% ====================================================================
\section{Zwierciadła płaskie i sferyczne}
% ====================================================================

\subsection{Zwierciadło płaskie}

Zwierciadło płaskie tworzy obraz \textbf{pozorny}, \textbf{prosty} i tej samej
wielkości co przedmiot.  Obraz leży symetrycznie za zwierciadłem.

\begin{itemize}
  \item Odległość obrazu za zwierciadłem = odległość przedmiotu przed zwierciadłem.
  \item Obraz jest \emph{lustrzanym odbiciem} (lateralne odwrócenie).
\end{itemize}

\subsubsection{Konstrukcja graficzna -- zwierciadło płaskie}

\begin{figure}[H]
\centering
\begin{tikzpicture}[scale=0.9]
  % Mirror
  \draw[thick] (0,-2.5) -- (0,2.5);
  \foreach \y in {-2.2,-1.8,...,2.2}
    \draw[gray!60] (0,\y) -- (-0.25,\y-0.25);
  % Object
  \draw[blue,thick,->] (2,0) -- (2,1.5) node[right]{$A$};
  \filldraw[blue] (2,0) circle (2pt) node[below right]{$B$};
  % Image (dashed)
  \draw[blue!40,dashed,->] (-2,0) -- (-2,1.5) node[left]{$A'$};
  \filldraw[blue!40] (-2,0) circle (2pt) node[below left]{$B'$};
  % Rays from A
  \draw[orange,->] (2,1.5) -- (0,1.5) -- (-0.5,2.8);
  \draw[orange,dashed] (0,1.5) -- (-2,1.5);
  \draw[orange,->] (2,1.5) -- (0,0.5) -- (-1.5,0.5);
  \draw[orange,dashed] (0,0.5) -- (-2,1.5);
  % Labels
  \node[below] at (1,0) {$d$};
  \node[below] at (-1,0) {$d'$};
  \draw[<->] (0,-0.3) -- (2,-0.3) node[midway,below]{$d$};
  \draw[<->] (0,-0.3) -- (-2,-0.3) node[midway,below]{$d'=d$};
\end{tikzpicture}
\caption{Powstawanie obrazu w zwierciadle płaskim. Przerywane linie -- przedłużenia
         promieni odbitych (obraz pozorny).}
\end{figure}

\subsection{Zwierciadło sferyczne}

\subsubsection{Ognisko i ogniskowa}

Zwierciadło sferyczne ma \textbf{środek krzywizny} $C$ i \textbf{ognisko} $F$
(dla zwierciadła wklęsłego -- rzeczywiste; dla wypukłego -- pozorne).

\begin{wzorbox}
\begin{equation}
  f = \frac{R}{2}
  \label{eq:ogniskowa_zwierciadla}
\end{equation}
\end{wzorbox}

\begin{objasnbox}
\textbf{Objaśnienie symboli:}
\begin{itemize}[leftmargin=2em]
  \oznaczenie{f}{Ogniskowa zwierciadła \quad [\textbf{m}]}
  \oznaczenie{R}{Promień krzywizny zwierciadła \quad [\textbf{m}]}
\end{itemize}
\textbf{Konwencja znaków:} ogniskowa wklęsłego $f>0$, wypukłego $f<0$.
\end{objasnbox}

\subsubsection{Równanie zwierciadła}

\begin{wzorbox}
\begin{equation}
  \frac{1}{f} = \frac{1}{d_o} + \frac{1}{d_i}
  \label{eq:zwierciadlo}
\end{equation}
\end{wzorbox}

\begin{objasnbox}
\textbf{Objaśnienie symboli:}
\begin{itemize}[leftmargin=2em]
  \oznaczenie{f}{Ogniskowa zwierciadła \quad [\textbf{m}]}
  \oznaczenie{d_o}{Odległość \textbf{przedmiotu} od zwierciadła \quad [\textbf{m}]}
  \oznaczenie{d_i}{Odległość \textbf{obrazu} od zwierciadła \quad [\textbf{m}]}
\end{itemize}
\textbf{Konwencja znaków (kartezjańska):}
\begin{itemize}
  \item $d_o > 0$ -- przedmiot przed zwierciadłem (rzeczywisty)
  \item $d_i > 0$ -- obraz przed zwierciadłem (rzeczywisty)
  \item $d_i < 0$ -- obraz za zwierciadłem (pozorny)
\end{itemize}
\end{objasnbox}

\subsubsection{Powiększenie obrazu}

\begin{wzorbox}
\begin{equation}
  m = -\frac{d_i}{d_o} = \frac{h_i}{h_o}
  \label{eq:powiekszenie_zwierciadlo}
\end{equation}
\end{wzorbox}

\begin{objasnbox}
\textbf{Objaśnienie symboli:}
\begin{itemize}[leftmargin=2em]
  \oznaczenie{m}{Powiększenie liniowe (bezwymiarowe)}
  \oznaczenie{d_i}{Odległość obrazu od zwierciadła \quad [\textbf{m}]}
  \oznaczenie{d_o}{Odległość przedmiotu od zwierciadła \quad [\textbf{m}]}
  \oznaczenie{h_i}{Wysokość obrazu \quad [\textbf{m}]}
  \oznaczenie{h_o}{Wysokość przedmiotu \quad [\textbf{m}]}
\end{itemize}
\textbf{Interpretacja:}
\begin{itemize}
  \item $|m|>1$ -- obraz powiększony
  \item $|m|<1$ -- obraz pomniejszony
  \item $m>0$   -- obraz prosty
  \item $m<0$   -- obraz odwrócony
\end{itemize}
\end{objasnbox}

\subsubsection{Obrazy w zwierciadle wklęsłym -- podsumowanie}

\begin{table}[H]
\centering
\renewcommand{\arraystretch}{1.4}
\begin{tabular}{lll}
\toprule
\textbf{Położenie przedmiotu} & \textbf{Charakter obrazu} & \textbf{Powiększenie} \\
\midrule
$d_o > 2f$ & rzeczywisty, odwrócony & $|m|<1$ \\
$d_o = 2f$ & rzeczywisty, odwrócony & $|m|=1$ \\
$f < d_o < 2f$ & rzeczywisty, odwrócony & $|m|>1$ \\
$d_o = f$ & obraz w nieskończoności & --- \\
$d_o < f$ & pozorny, prosty & $|m|>1$ \\
\bottomrule
\end{tabular}
\caption{Obrazy w zwierciadle wklęsłym}
\end{table}

\subsubsection{Konstrukcja graficzna -- zwierciadło sferyczne}

Stosuje się trzy \textbf{promienie pomocnicze}:
\begin{enumerate}
  \item Promień równoległy do osi optycznej $\to$ po odbiciu przechodzi przez ognisko $F$.
  \item Promień przez ognisko $F$ $\to$ po odbiciu biegnie równolegle do osi.
  \item Promień przez środek krzywizny $C$ $\to$ odbija się wzdłuż siebie.
\end{enumerate}

\begin{figure}[H]
\centering
\begin{tikzpicture}[scale=1.0]
  % Axis
  \draw[->] (-5,0) -- (2,0) node[right]{oś opt.};
  % Mirror (parabolic approximation)
  \draw[thick,blue!70] plot[domain=-2:2,samples=40] ({-0.15*\x*\x},\x);
  % Focal point and centre
  \filldraw[red] (-1,0) circle (2pt) node[below]{$F$};
  \filldraw[red] (-2,0) circle (2pt) node[below]{$C$};
  \node[right] at (0,0) {$V$};
  % Object
  \draw[thick,green!60!black,->] (-3.5,0) -- (-3.5,1.2) node[above]{$P$};
  % Ray 1: parallel to axis -> through F
  \draw[orange,->] (-3.5,1.2) -- (0,1.2);
  \draw[orange,->] (0,1.2) -- (-1.8,-1.0);
  % Ray 2: through C -> reflects back
  \draw[purple,->] (-3.5,1.2) -- (-0.06,0.2);
  \draw[purple,->] (-0.06,0.2) -- (-3.5,1.2) +(-3.5+0.06, 1.2-0.2);
  % Better: ray through F then parallel
  \draw[cyan,->] (-3.5,1.2) -- (-0.04,0.65);
  \draw[cyan,->] (-0.04,0.65) -- (-2.5,-0.8);
  % Image label
  \node[below right, green!60!black] at (-2,-0.35) {obraz $P'$};
\end{tikzpicture}
\caption{Szkic konstrukcji obrazu w zwierciadle wklęsłym (schemat poglądowy).}
\end{figure}

\newpage
% ====================================================================
\section{Soczewki}
% ====================================================================

\subsection{Soczewki skupiające i rozpraszające}

\begin{itemize}
  \item \textbf{Soczewka skupiająca (wypukła)} -- grubsza w środku; ognisko
        \emph{rzeczywiste} (po stronie obrazu); $f>0$.
  \item \textbf{Soczewka rozpraszająca (wklęsła)} -- grubsza na brzegach; ognisko
        \emph{pozorne} (po tej samej stronie co przedmiot); $f<0$.
\end{itemize}

\subsubsection{Ognisko i ogniskowa soczewki}

Ogniskowa $f$ soczewki skupiającej to odległość, w jakiej zbiega się wiązka promieni
równoległych do osi optycznej po przejściu przez soczewkę.

\subsubsection{Równanie soczewki cienkiej}

\begin{wzorbox}
\begin{equation}
  \frac{1}{f} = \frac{1}{d_o} + \frac{1}{d_i}
  \label{eq:soczewka_cienka}
\end{equation}
\end{wzorbox}

\begin{objasnbox}
\textbf{Objaśnienie symboli:}
\begin{itemize}[leftmargin=2em]
  \oznaczenie{f}{Ogniskowa soczewki \quad [\textbf{m}]; dla skupiającej $f>0$,
                dla rozpraszającej $f<0$}
  \oznaczenie{d_o}{Odległość przedmiotu od soczewki \quad [\textbf{m}]; $d_o>0$
                po lewej stronie soczewki}
  \oznaczenie{d_i}{Odległość obrazu od soczewki \quad [\textbf{m}]; $d_i>0$
                obraz rzeczywisty (po prawej), $d_i<0$ obraz pozorny (po lewej)}
\end{itemize}
\end{objasnbox}

\subsubsection{Zdolność skupiająca soczewki}

Często posługujemy się zdolnością skupiającą (mocą) soczewki:
\begin{wzorbox}
\begin{equation}
  D = \frac{1}{f}
  \label{eq:zdolnosc_skupiajaca}
\end{equation}
\end{wzorbox}

\begin{objasnbox}
\textbf{Objaśnienie symboli:}
\begin{itemize}[leftmargin=2em]
  \oznaczenie{D}{Zdolność skupiająca \quad [\textbf{dioptria, D}] $= \si{m^{-1}}$}
  \oznaczenie{f}{Ogniskowa \quad [\textbf{m}]}
\end{itemize}
\end{objasnbox}

\subsubsection{Równanie szlifierzy soczewek (soczewki rzeczywiste)}

Dla soczewki z materiału o współczynniku załamania $n_L$ zanurzonej
w ośrodku o współczynniku $n_0$:

\begin{wzorbox}
\begin{equation}
  \frac{1}{f} = \left(\frac{n_L}{n_0} - 1\right)
                \left(\frac{1}{R_1} - \frac{1}{R_2}\right)
  \label{eq:szlifierzy}
\end{equation}
\end{wzorbox}

\begin{objasnbox}
\textbf{Objaśnienie symboli:}
\begin{itemize}[leftmargin=2em]
  \oznaczenie{f}{Ogniskowa soczewki \quad [\textbf{m}]}
  \oznaczenie{n_L}{Współczynnik załamania materiału soczewki (bezwymiarowy)}
  \oznaczenie{n_0}{Współczynnik załamania ośrodka otaczającego (dla powietrza $n_0\approx1$)}
  \oznaczenie{R_1}{Promień krzywizny pierwszej powierzchni soczewki \quad [\textbf{m}];\\
               $R_1>0$ -- środek krzywizny po prawej stronie}
  \oznaczenie{R_2}{Promień krzywizny drugiej powierzchni soczewki \quad [\textbf{m}];\\
               $R_2<0$ -- środek krzywizny po lewej stronie (typowe dla bikonweksowej)}
\end{itemize}
\end{objasnbox}

Dla powietrza ($n_0=1$) wzór upraszcza się do:
\begin{equation}
  \frac{1}{f} = (n_L - 1)\left(\frac{1}{R_1} - \frac{1}{R_2}\right)
\end{equation}

\subsubsection{Powiększenie przez soczewkę}

\begin{wzorbox}
\begin{equation}
  m = \frac{h_i}{h_o} = -\frac{d_i}{d_o}
  \label{eq:powiekszenie_soczewka}
\end{equation}
\end{wzorbox}

\begin{objasnbox}
\textbf{Objaśnienie symboli:}
\begin{itemize}[leftmargin=2em]
  \oznaczenie{m}{Powiększenie liniowe (bezwymiarowe)}
  \oznaczenie{h_i}{Rozmiar obrazu \quad [\textbf{m}]}
  \oznaczenie{h_o}{Rozmiar przedmiotu \quad [\textbf{m}]}
  \oznaczenie{d_i}{Odległość obrazu \quad [\textbf{m}]}
  \oznaczenie{d_o}{Odległość przedmiotu \quad [\textbf{m}]}
\end{itemize}
\end{objasnbox}

\subsubsection{Obrazy w soczewce skupiającej -- podsumowanie}

\begin{table}[H]
\centering
\renewcommand{\arraystretch}{1.4}
\begin{tabular}{lll}
\toprule
\textbf{Położenie przedmiotu} & \textbf{Charakter obrazu} & \textbf{Powiększenie} \\
\midrule
$d_o > 2f$        & rzeczywisty, odwrócony & $|m|<1$ \\
$d_o = 2f$        & rzeczywisty, odwrócony & $|m|=1$ \\
$f < d_o < 2f$    & rzeczywisty, odwrócony & $|m|>1$ \\
$d_o = f$         & brak obrazu (wiązka równoległa) & --- \\
$d_o < f$         & pozorny, prosty       & $|m|>1$ \\
\bottomrule
\end{tabular}
\caption{Obrazy w soczewce skupiającej}
\end{table}

\subsubsection{Konstrukcja geometryczna obrazów w soczewce}

Podobnie jak dla zwierciadeł stosuje się trzy promienie pomocnicze:
\begin{enumerate}
  \item Promień równoległy do osi optycznej $\to$ po przejściu przez soczewkę skupiającą
        przechodzi przez ognisko $F_2$ po prawej stronie.
  \item Promień przez środek optyczny soczewki $O$ $\to$ nie zmienia kierunku.
  \item Promień przez ognisko $F_1$ po lewej $\to$ po przejściu wychodzi
        równolegle do osi optycznej.
\end{enumerate}

\newpage
% ====================================================================
\section{Odbicie i załamanie światła}
% ====================================================================

\subsection{Współczynnik załamania światła}

Gdy światło przechodzi z próżni (lub powietrza) do ośrodka optycznie gęstszego,
jego prędkość maleje.  Współczynnik załamania opisuje tę relację:

\begin{wzorbox}
\begin{equation}
  n = \frac{c}{v}
  \label{eq:wspolczynnik_zalamaniia}
\end{equation}
\end{wzorbox}

\begin{objasnbox}
\textbf{Objaśnienie symboli:}
\begin{itemize}[leftmargin=2em]
  \oznaczenie{n}{Bezwzględny współczynnik załamania (bezwymiarowy; $n \geq 1$)}
  \oznaczenie{c}{Prędkość światła w próżni, $c \approx 3\times10^8\,\si{m/s}$}
  \oznaczenie{v}{Prędkość fazowa fali w danym ośrodku \quad [\textbf{m/s}]}
\end{itemize}
\textbf{Przykłady:}
\begin{itemize}
  \item Powietrze: $n \approx 1{,}000$
  \item Woda: $n \approx 1{,}33$
  \item Szkło optyczne: $n \approx 1{,}5$
  \item Diament: $n \approx 2{,}42$
\end{itemize}
\end{objasnbox}

\subsection{Kąty padania, odbicia i załamania}

\begin{itemize}
  \item \textbf{Kąt padania} $\theta_1$ -- kąt między promieniem \emph{padającym}
        a normalną do powierzchni granicznej w punkcie padania.
  \item \textbf{Kąt odbicia} $\theta_r$ -- kąt między promieniem \emph{odbitym}
        a normalną.
  \item \textbf{Kąt załamania} $\theta_2$ -- kąt między promieniem \emph{załamanym}
        a normalną po drugiej stronie powierzchni.
\end{itemize}

Wszystkie kąty mierzone są od \textbf{normalnej} do powierzchni, nie od samej powierzchni.

\subsection{Prawo odbicia}

\begin{wzorbox}
\begin{equation}
  \theta_r = \theta_1
  \label{eq:prawo_odbicia}
\end{equation}
\end{wzorbox}

\begin{objasnbox}
\textbf{Objaśnienie:}
\begin{itemize}[leftmargin=2em]
  \oznaczenie{\theta_1}{Kąt padania (od normalnej) \quad [°] lub [\textbf{rad}]}
  \oznaczenie{\theta_r}{Kąt odbicia (od normalnej) \quad [°] lub [\textbf{rad}]}
\end{itemize}
Promień padający, odbity i normalna leżą w jednej płaszczyźnie.
\end{objasnbox}

\subsection{Prawo Snella (Snelliusa)}

\begin{wzorbox}
\begin{equation}
  n_1 \sin\theta_1 = n_2 \sin\theta_2
  \label{eq:snell}
\end{equation}
\end{wzorbox}

\begin{objasnbox}
\textbf{Objaśnienie symboli:}
\begin{itemize}[leftmargin=2em]
  \oznaczenie{n_1}{Współczynnik załamania ośrodka, w którym biegnie promień padający
                   (bezwymiarowy)}
  \oznaczenie{\theta_1}{Kąt padania \quad [°] lub [\textbf{rad}]}
  \oznaczenie{n_2}{Współczynnik załamania ośrodka, w którym biegnie promień załamany
                   (bezwymiarowy)}
  \oznaczenie{\theta_2}{Kąt załamania \quad [°] lub [\textbf{rad}]}
\end{itemize}
\end{objasnbox}

\subsubsection{Całkowite wewnętrzne odbicie}

Gdy $n_1 > n_2$ i kąt padania $\theta_1$ przekracza \textbf{kąt graniczny} $\theta_c$,
następuje całkowite wewnętrzne odbicie (brak promienia załamanego):

\begin{wzorbox}
\begin{equation}
  \sin\theta_c = \frac{n_2}{n_1}
  \label{eq:kat_graniczny}
\end{equation}
\end{wzorbox}

\begin{objasnbox}
\textbf{Objaśnienie symboli:}
\begin{itemize}[leftmargin=2em]
  \oznaczenie{\theta_c}{Kąt graniczny (krytyczny) całkowitego wewnętrznego odbicia
                        \quad [°] lub [\textbf{rad}]}
  \oznaczenie{n_1}{Współczynnik załamania ośrodka gęstszego optycznie ($n_1>n_2$)}
  \oznaczenie{n_2}{Współczynnik załamania ośrodka rzadszego optycznie}
\end{itemize}
\textbf{Zastosowanie:} światłowody, pryzmaty odchylające.
\end{objasnbox}

\begin{figure}[H]
\centering
\begin{tikzpicture}[scale=1.1]
  % Interface
  \fill[blue!8] (-3,-2) rectangle (3,0);
  \draw[thick] (-3,0) -- (3,0);
  \node[right] at (3,0) {granica ośrodków};
  \node[above right] at (-3,0.1) {$n_1$};
  \node[below right] at (-3,-0.1) {$n_2 > n_1$};
  % Normal
  \draw[dashed] (0,-2) -- (0,2);
  % Incident ray
  \draw[orange,very thick,->] (-2,1.6) -- (0,0) node[midway,above]{padający};
  % Reflected ray
  \draw[orange!70,very thick,->] (0,0) -- (2,1.6) node[midway,above]{odbity};
  % Refracted ray
  \draw[red,very thick,->] (0,0) -- (0.8,-1.8) node[midway,right]{załamany};
  % Angles
  \draw[<->,gray] (0,0.8) arc (90:128:0.8) node[midway,above left,scale=0.8]{$\theta_1$};
  \draw[<->,gray] (0,0.8) arc (90:52:0.8) node[midway,above right,scale=0.8]{$\theta_r$};
  \draw[<->,gray] (0,-0.8) arc (270:295:0.8) node[midway,below right,scale=0.8]{$\theta_2$};
\end{tikzpicture}
\caption{Odbicie i załamanie światła na granicy ośrodków.}
\end{figure}

\newpage
% ====================================================================
\section{Rozszczepienie światła}
% ====================================================================

\subsection{Dyspersja chromatyczna}

Dyspersja chromatyczna to zjawisko zależności współczynnika załamania $n$ od
długości fali $\lambda$ (lub częstotliwości $\nu$) światła.  W ośrodku dyspersyjnym:

\begin{wzorbox}
\begin{equation}
  n = n(\lambda)
  \label{eq:dyspersja}
\end{equation}
\end{wzorbox}

Zwykle w szkle $n$ maleje wraz ze wzrostem $\lambda$ -- jest to
\textbf{dyspersja normalna}.  Współczynnik dyspersji Abbego opisuje stopień
dyspersji materiału:

\begin{wzorbox}
\begin{equation}
  V_d = \frac{n_d - 1}{n_F - n_C}
  \label{eq:liczba_abbego}
\end{equation}
\end{wzorbox}

\begin{objasnbox}
\textbf{Objaśnienie symboli:}
\begin{itemize}[leftmargin=2em]
  \oznaczenie{V_d}{Liczba Abbego (współczynnik dyspersji) -- bezwymiarowa;
                   większa wartość = mniejsza dyspersja}
  \oznaczenie{n_d}{Współczynnik załamania dla linii $d$ helu $\lambda_d = 587{,}6\,\si{nm}$
                   (żółta)}
  \oznaczenie{n_F}{Współczynnik załamania dla linii $F$ wodoru $\lambda_F = 486{,}1\,\si{nm}$
                   (niebieska)}
  \oznaczenie{n_C}{Współczynnik załamania dla linii $C$ wodoru $\lambda_C = 656{,}3\,\si{nm}$
                   (czerwona)}
\end{itemize}
\end{objasnbox}

\subsection{Widmo promieniowania elektromagnetycznego}

\begin{table}[H]
\centering
\renewcommand{\arraystretch}{1.4}
\begin{tabular}{llll}
\toprule
\textbf{Zakres} & \textbf{Skrót} &
\textbf{Długość fali $\lambda$} & \textbf{Częstotliwość $\nu$} \\
\midrule
Nadfiolet (ultrafiolet) & UV  & $100 - 400\,\si{nm}$   & $750\,\si{THz} - 30\,\si{PHz}$ \\
Widzialne              & VIS  & $380 - 780\,\si{nm}$   & $385 - 789\,\si{THz}$ \\
\quad fiolet           &      & $380 - 450\,\si{nm}$   & $670 - 789\,\si{THz}$ \\
\quad niebieski        &      & $450 - 490\,\si{nm}$   & $612 - 670\,\si{THz}$ \\
\quad zielony          &      & $490 - 570\,\si{nm}$   & $526 - 612\,\si{THz}$ \\
\quad żółty            &      & $570 - 590\,\si{nm}$   & $508 - 526\,\si{THz}$ \\
\quad pomarańczowy     &      & $590 - 620\,\si{nm}$   & $484 - 508\,\si{THz}$ \\
\quad czerwony         &      & $620 - 780\,\si{nm}$   & $385 - 484\,\si{THz}$ \\
Podczerwień (infrared) & IR   & $780\,\si{nm} - 1\,\si{mm}$ & $300\,\si{GHz} - 385\,\si{THz}$ \\
\bottomrule
\end{tabular}
\caption{Zakresy promieniowania elektromagnetycznego}
\end{table}

Zależność długości fali i częstotliwości:
\begin{wzorbox}
\begin{equation}
  c = \lambda \nu \quad \Longrightarrow \quad \nu = \frac{c}{\lambda}
  \label{eq:predkosc_fali}
\end{equation}
\end{wzorbox}

\begin{objasnbox}
\textbf{Objaśnienie symboli:}
\begin{itemize}[leftmargin=2em]
  \oznaczenie{c}{Prędkość światła w próżni, $c \approx 3\times10^8\,\si{m/s}}
  \oznaczenie{\lambda}{Długość fali \quad [\textbf{m}]; praktycznie [\textbf{nm}]
                       (1\,nm $= 10^{-9}$\,m)}
  \oznaczenie{\nu}{Częstotliwość \quad [\textbf{Hz}]; praktycznie [\textbf{THz}]
                   (1\,THz $= 10^{12}$\,Hz)}
\end{itemize}
\end{objasnbox}

\subsection{Bieg promienia w pryzmacie}

\begin{figure}[H]
\centering
\begin{tikzpicture}[scale=1.2]
  % Prism
  \draw[thick,fill=blue!10] (0,0) -- (3,0) -- (1.5,2.6) -- cycle;
  % Apex angle label
  \node at (1.5,2.0) {$\varphi$};
  % Incident ray
  \draw[orange,very thick,->] (-1,1.8) -- (0.6,1.3)
        node[midway,above left]{promień padający};
  % Ray inside prism
  \draw[orange!70,thick] (0.6,1.3) -- (2.2,0.9);
  % Refracted ray
  \draw[red,very thick,->] (2.2,0.9) -- (3.8,0.3)
        node[midway,below right]{promień wychodzący};
  % Deviation angle
  \draw[dashed,gray] (-1,1.8) -- (4.0,1.8);
  \draw[<->,blue!80] (3.0,1.8) -- (3.8,0.3) node[midway,right]{$\delta$};
  % Apex
  \node[above] at (1.5,2.6) {$A$};
\end{tikzpicture}
\caption{Bieg promienia w pryzmacie. $\varphi$ -- kąt łamiący, $\delta$ -- kąt odchylenia.}
\end{figure}

\subsubsection{Kąt łamiący pryzmatu i odchylenie promienia}

\begin{wzorbox}
\begin{equation}
  \delta = (\theta_1 + \theta_2') - \varphi
  \label{eq:odchylenie_pryzmatu}
\end{equation}
\end{wzorbox}

\begin{objasnbox}
\textbf{Objaśnienie symboli:}
\begin{itemize}[leftmargin=2em]
  \oznaczenie{\delta}{Kąt odchylenia promienia przez pryzmat \quad [°] lub [\textbf{rad}]}
  \oznaczenie{\theta_1}{Kąt padania na pierwszą powierzchnię pryzmatu \quad [°]}
  \oznaczenie{\theta_2'}{Kąt wyjścia z drugiej powierzchni pryzmatu \quad [°]}
  \oznaczenie{\varphi}{Kąt łamiący (wierzchołkowy) pryzmatu \quad [°]}
\end{itemize}
\end{objasnbox}

Na wewnątrz pryzmatu prawo Snella stosuje się dwukrotnie:
\begin{equation}
  n_0 \sin\theta_1 = n_p \sin\theta_1'
  \qquad\text{(pierwsza powierzchnia)}
\end{equation}
\begin{equation}
  n_p \sin\theta_2 = n_0 \sin\theta_2'
  \qquad\text{(druga powierzchnia)}
\end{equation}

gdzie kąty wewnątrz pryzmatu spełniają:
\begin{equation}
  \theta_1' + \theta_2 = \varphi
\end{equation}

\begin{objasnbox}
\textbf{Objaśnienie symboli:}
\begin{itemize}[leftmargin=2em]
  \oznaczenie{n_0}{Współczynnik załamania ośrodka (powietrze $\approx 1$)}
  \oznaczenie{n_p}{Współczynnik załamania materiału pryzmatu}
  \oznaczenie{\theta_1'}{Kąt załamania na pierwszej powierzchni (wewnątrz pryzmatu)}
  \oznaczenie{\theta_2}{Kąt padania na drugą powierzchnię (wewnątrz pryzmatu)}
\end{itemize}
\end{objasnbox}

\subsubsection{Najmniejszy kąt odchylenia $\delta_{\min}$}

Minimalne odchylenie zachodzi przy symetrycznym biegu promienia przez pryzmat
(promień wewnątrz pryzmatu jest równoległy do podstawy).

\begin{wzorbox}
\begin{equation}
  n_p = \frac{\sin\!\left(\dfrac{\varphi + \delta_{\min}}{2}\right)}
             {\sin\!\left(\dfrac{\varphi}{2}\right)}
  \label{eq:min_odchylenie}
\end{equation}
\end{wzorbox}

\begin{objasnbox}
\textbf{Objaśnienie symboli:}
\begin{itemize}[leftmargin=2em]
  \oznaczenie{n_p}{Współczynnik załamania materiału pryzmatu (bezwymiarowy)}
  \oznaczenie{\varphi}{Kąt łamiący pryzmatu \quad [°]}
  \oznaczenie{\delta_{\min}}{Minimalny kąt odchylenia promienia \quad [°]}
\end{itemize}
Wzór ten służy do wyznaczania $n_p$ na podstawie pomiarów $\varphi$ i $\delta_{\min}$.
\end{objasnbox}

Przy małych kątach łamiących ($\varphi \ll 1$) wzór upraszcza się:
\begin{equation}
  \delta \approx (n_p - 1)\,\varphi
  \label{eq:pryzmnat_male_katy}
\end{equation}

% ====================================================================
\subsection{Zadanie: pryzmat równoboczny w ośrodku o $n_2 = 1{,}33$}
% ====================================================================

\textbf{Treść:} Pryzmat wykonany jest z materiału o współczynniku załamania
$n_1 = 1$, a jego przekrój stanowi trójkąt równoboczny ($\varphi = 60°$).
Współczynnik załamania otoczenia wynosi $n_2 = 1{,}33$.
Promień pada prostopadle na jedną z powierzchni pryzmatu.
Narysuj bieg promienia i oblicz kąt odchylenia $\delta$.

\begin{figure}[H]
\centering
\begin{tikzpicture}[scale=1.3]
  % --- equilateral triangle: A=(0,0), B=(3,0), C=(1.5,2.598) ---
  \draw[thick, fill=blue!10] (0,0) -- (3,0) -- (1.5,2.598) -- cycle;
  \node[left]  at (0,0)      {$A$};
  \node[right] at (3,0)      {$B$};
  \node[above] at (1.5,2.598){$C$};
  \node at (1.5,0.35) {$\varphi = 60°$};

  % --- entry point P on AC (3/4 from A), exit point Q on BC (midpoint) ---
  % P = (1.125, 1.949),  Q = (2.25, 1.299)
  % inside-prism direction: (sqrt(3)/2, -1/2) = (0.866, -0.5)

  % --- normal to AC at P (direction (0.866,-0.5), draw through P) ---
  \draw[dashed, gray] (1.125-1.1*0.866, 1.949+1.1*0.5)
                   -- (1.125+1.3*0.866, 1.949-1.3*0.5);
  % that is from (0.173, 2.499) to (2.252, 1.299) approx

  % --- incoming ray perpendicular to AC: arrives from outside along (0.866,-0.5) ---
  \draw[orange, very thick, ->]
        (1.125-1.5*0.866, 1.949+1.5*0.5) -- (1.125, 1.949)
        node[midway, above left] {promień padający ($\theta_1=0°$)};

  % --- right-angle mark at P (ray perpendicular to surface) ---
  % surface AC direction: (0.5, 0.866); normal: (0.866,-0.5)
  % small square: P + 0.18*(0.5,0.866) + 0.18*(0.866,-0.5)
  \draw (1.125+0.09+0.156, 1.949+0.156-0.09)
     -- (1.125+0.156,      1.949-0.09)
     -- (1.125,            1.949);
  \node[left, font=\footnotesize] at (1.0, 1.949) {$\theta_1=0°$};

  % --- ray inside prism: P -> Q ---
  \draw[orange!70, thick] (1.125, 1.949) -- (2.25, 1.299);

  % --- normal to BC at Q (outward direction (0.866,0.5), draw through Q) ---
  \draw[dashed, gray] (2.25-1.4*0.866, 1.299-1.4*0.5)
                   -- (2.25+1.3*0.866, 1.299+1.3*0.5);
  % from (1.038, 0.599) to (3.376, 1.949)

  % --- angle theta_2 = 60 deg at Q (between inside ray and normal) ---
  \node[above left, font=\footnotesize] at (2.10, 1.30) {$\theta_2{=}60°$};

  % --- exit ray from Q in direction (0.9828, -0.1849) ---
  \draw[red, very thick, ->]
        (2.25, 1.299) -- (2.25+2.0*0.9828, 1.299-2.0*0.1849)
        node[midway, below right] {promień wychodzący};

  % --- angle theta_2' = 40.6 deg at Q (between exit ray and normal) ---
  \node[right, font=\footnotesize] at (3.40, 1.95) {$\theta_2'{=}40{,}6°$};

  % --- extension of incident-direction from Q (dashed) to show deviation ---
  \draw[dashed, orange!60]
        (2.25, 1.299) -- (2.25+2.0*0.866, 1.299-2.0*0.5);
  % to (3.982, 0.299)

  % --- deviation angle arc delta at Q ---
  \draw[<->, blue!80, thick]
        (2.25+1.8*0.9828, 1.299-1.8*0.1849) arc[start angle=-10.7,
        end angle=-30, radius=1.8]
        node[midway, right, font=\footnotesize] {$|\delta|\approx19{,}4°$};

  % --- apex angle label ---
  \node[below, font=\footnotesize] at (1.5, 2.45) {$60°$};
\end{tikzpicture}
\caption{Bieg promienia przez pryzmat równoboczny ($n_1=1$) zanurzony
         w ośrodku o $n_2=1{,}33$.
         Promień wchodzi prostopadle przez lewą ścianę (brak załamania),
         odbija się od normalnej na prawej ściance ku wierzchołkowi.}
\end{figure}

\subsubsection*{Rozwiązanie krok po kroku}

\paragraph{Krok 1 – pierwsza powierzchnia (wejście).}
Promień pada prostopadle na ścianę, zatem kąt padania wynosi $\theta_1 = 0°$.
Z prawa Snella:
\[
  n_2 \sin\theta_1 = n_1 \sin\theta_1'
  \;\Longrightarrow\;
  1{,}33 \cdot \sin 0° = 1 \cdot \sin\theta_1'
  \;\Longrightarrow\;
  \theta_1' = 0°.
\]
Brak załamania -- promień przechodzi bez odchylenia.

\paragraph{Krok 2 – geometria wewnątrz pryzmatu.}
Dla trójkąta równobocznego kąty wewnętrzne pryzmatu spełniają:
\[
  \theta_1' + \theta_2 = \varphi
  \;\Longrightarrow\;
  \theta_2 = \varphi - \theta_1' = 60° - 0° = 60°.
\]
Promień uderza w drugą ścianę pod kątem $\theta_2 = 60°$.

\paragraph{Krok 3 – druga powierzchnia (wyjście).}
Promień przechodzi z ośrodka o $n_1 = 1$ do otoczenia o $n_2 = 1{,}33$
(ze rzadszego do gęstszego -- zagięcie ku normalnej):
\[
  n_1 \sin\theta_2 = n_2 \sin\theta_2'
  \;\Longrightarrow\;
  1 \cdot \sin 60° = 1{,}33 \cdot \sin\theta_2'
\]
\[
  \sin\theta_2' = \frac{\sin 60°}{1{,}33}
               = \frac{\sqrt{3}/2}{1{,}33}
               \approx \frac{0{,}8660}{1{,}33}
               \approx 0{,}6511
\]
\[
  \boxed{\theta_2' = \arcsin(0{,}6511) \approx 40{,}6°}
\]

\paragraph{Krok 4 – całkowite odchylenie.}

\begin{wzorbox}
\[
  \delta = (\theta_1 + \theta_2') - \varphi
         = (0° + 40{,}6°) - 60°
         = -19{,}4°
\]
\end{wzorbox}

\begin{objasnbox}
Ujemna wartość $\delta$ oznacza, że promień odchyla się
\textbf{w kierunku wierzchołka} pryzmatu (a nie ku podstawie jak
w klasycznym pryzmacie szklanym w powietrzu).
Wynika to z faktu, że $n_1 < n_2$: na wyjściu promień ugina się
\emph{ku normalnej}, a nie od niej.
Kąt odchylenia wynosi więc $|\delta| \approx \mathbf{19{,}4°}$.
\end{objasnbox}

% ====================================================================
\section{Zestawienie wszystkich wzorów}
% ====================================================================

\begin{table}[H]
\centering
\renewcommand{\arraystretch}{1.6}
\begin{tabular}{lll}
\toprule
\textbf{Temat} & \textbf{Wzór} & \textbf{Nr równania} \\
\midrule
Fotometria -- światłość     & $I = d\Phi/d\Omega$ & \eqref{eq:swiatlosc} \\
Prawo odwrotnych kwadratów  & $E = I/r^2$         & \eqref{eq:prawo_odwrotnych_kwadratow} \\
Oświetlenie (kąt)           & $E = I\cos\theta/r^2$ & \eqref{eq:lambert} \\
Ogniskowa zwierciadła       & $f = R/2$           & \eqref{eq:ogniskowa_zwierciadla} \\
Równanie zwierciadła        & $1/f = 1/d_o + 1/d_i$ & \eqref{eq:zwierciadlo} \\
Powiększenie                & $m = -d_i/d_o$      & \eqref{eq:powiekszenie_zwierciadlo} \\
Równanie soczewki cienkiej  & $1/f = 1/d_o + 1/d_i$ & \eqref{eq:soczewka_cienka} \\
Równanie szlifierzy soczewek & $1/f=(n_L/n_0-1)(1/R_1-1/R_2)$ & \eqref{eq:szlifierzy} \\
Zdolność skupiająca         & $D = 1/f$           & \eqref{eq:zdolnosc_skupiajaca} \\
Wsp. załamania              & $n = c/v$           & \eqref{eq:wspolczynnik_zalamaniia} \\
Prawo odbicia               & $\theta_r = \theta_1$ & \eqref{eq:prawo_odbicia} \\
Prawo Snella                & $n_1\sin\theta_1 = n_2\sin\theta_2$ & \eqref{eq:snell} \\
Kąt graniczny               & $\sin\theta_c = n_2/n_1$ & \eqref{eq:kat_graniczny} \\
Prędkość fali               & $c = \lambda\nu$    & \eqref{eq:predkosc_fali} \\
Odchylenie w pryzmacie      & $\delta = (\theta_1+\theta_2')-\varphi$ & \eqref{eq:odchylenie_pryzmatu} \\
Min. odchylenie w pryzmacie & $n_p = \sin((\varphi+\delta_{\min})/2)/\sin(\varphi/2)$ & \eqref{eq:min_odchylenie} \\
\bottomrule
\end{tabular}
\caption{Zestawienie wzorów kolokwialnych}
\end{table}

\end{document}
